\textbf{COLUNI 2025 -- Questão 01}\par\medskip Leia o texto abaixo e responda as questões de 1 a 3.\par\medskip

\begin{center}\textit{Meu Rosário}\end{center}

\noindent
\begin{minipage}[t]{0.49\linewidth}
Meu rosário é feito de contas negras e mágicas.\\
Nas contas de meu rosário eu canto Mamãe Oxum e falo\\
padres-nossos, ave-marias.\\
Do meu rosário eu ouço os longínquos batuques do\\
meu povo\\
e encontro na memória mal-adormecida\\
as rezas dos meses de maio de minha infância.\\
As coroações da Senhora, onde as meninas negras,\\
apesar do desejo de coroar a Rainha,\\
tinham de se contentar em ficar ao pé do altar\\
lançando flores.\\
As contas do meu rosário fizeram calos\\
nas minhas mãos,\\
pois são contas do trabalho na terra, nas fábricas,\\
nas casas, nas escolas, nas ruas, no mundo.\\
As contas do meu rosário são contas vivas.\\
(Alguém disse que um dia a vida é uma oração,\\
eu diria porém que há vidas-blasfemas).\\
Nas contas de meu rosário eu teço entumecidos\\
sonhos de esperanças.
\end{minipage}\hfill
\begin{minipage}[t]{0.49\linewidth}
Nas contas do meu rosário eu vejo rostos escondidos\\
por visíveis e invisíveis grades\\
e embalo a dor da luta perdida nas contas\\
do meu rosário.\\
Nas contas de meu rosário eu canto, eu grito, eu calo.\\
Do meu rosário eu sinto o borbulhar da fome\\
No estômago, no coração e nas cabeças vazias.\\
Quando debulho as contas de meu rosário,\\
eu falo de mim mesma em outro nome.\\
E sonho nas contas de meu rosário lugares, pessoas,\\
vidas que pouco a pouco descubro reais.\\
Vou e volto por entre as contas de meu rosário,\\
que são pedras marcando-me o corpo-caminho.\\
E neste andar de contas-pedras,\\
o meu rosário se transmuda em tinta,\\
me guia o dedo,\\
me insinua a poesia.\\
E depois de macerar conta por conto do meu rosário,\\
me acho aqui eu mesma\\
e descubro que ainda me chamo Maria.
\end{minipage}

\begin{center}(EVARISTO, Conceição. \textit{Poemas da recordação e outros movimentos}. Rio de Janeiro: Malê, 2017. p. 38−39)\end{center}

Leia o excerto abaixo.

\par

O poeta afro-brasileiro é um partícipe ativo da presença negra na literatura brasileira, assim como possui plena consciência da necessidade de uma revisão crítica da História oficial que minimiza o passado de séculos de escravidão e a exclusão social que se perpetua para a maioria de nossos pares na contemporaneidade. (Disponível em: https://vermelho.org.br/2010/11/21/conceicao-evaristo-poemas-da-recordacao-e-outros-movimentos/. Acesso em: 01 jul. 2024)

\par

A partir do excerto e do poema \textit{Meu Rosário}, assinale a afirmativa INCORRETA.\begin{enumerate}
\item Há destaque para a ancestralidade marcada pela religião, bem como há referência à mãe.
\item Trata especialmente da descrição da dor, do sofrimento negro e da sua desesperança incisiva.
\item Há a presença das memórias do povo africano que vinham sendo adormecidas pelo povo colonizador.
\item Trata de papéis atribuídos aos afro-brasileiros para apontar a perspectiva de força e trabalho árduo.\end{enumerate}
